\section{Introduction}

We set out to measure how dataset documentation quality and benchmark concentration predict ML reproducibility failure --- and discovered that the metadata infrastructure we assumed existed does not. What began as a regression study became an infrastructure characterization: HuggingFace Hub and OpenML, the two most natural APIs for extracting dataset metadata at scale, achieve only 30\% and 22\% coverage respectively for the benchmark datasets used in Raff's 255-paper reproducibility corpus \citep{Raff2019Step}. The gap is not random. It is systematic, domain-specific, and --- once its root cause is understood --- immediately actionable.

The ML reproducibility crisis is well-documented. \citet{Raff2019Step} found that approximately 50\% of top-venue ML papers from 1984--2017 fail independent reproduction attempts. Subsequent work has proposed checklists \citep{Pineau2021Improving}, code sharing mandates, and dataset documentation frameworks \citep{Gebru2021Datasheets} as interventions. What is missing is an empirical linkage between \emph{metadata-observable signals} --- signals derivable from public APIs without re-running experiments --- and reproducibility failure outcomes. If such a linkage exists, it would enable automated pre-review flagging of high-risk papers without requiring the prohibitive cost of manual re-execution.

The deeper problem is methodological. We lack evidence that the documentation quality signal (HuggingFace dataset card field-presence, measuring how completely a dataset's scope, limitations, and intended use are specified) or the benchmark concentration signal (Herfindahl-Hirschman Index over pre-publication OpenML run counts, measuring how overfit-prone a dataset has become through concentrated community use) actually predict Raff's binary reproducibility labels. Two parallel literatures --- dataset documentation \citep{Gebru2021Datasheets, Bender2021Dangers} and ML reproducibility \citep{Raff2019Step, Pineau2021Improving} --- have advanced independently without an empirical bridge. The key question --- do metadata-observable misuse signals predict which papers fail reproduction? --- remains open.

The gap in existing work is structural: no study has tested this linkage because it requires querying post-2019 APIs (HuggingFace Hub launched 2019; OpenML's Python API matured around 2018--2020) against a pre-2018 corpus (Raff's 1984--2017 papers). We designed a systematic infrastructure feasibility study to determine whether the required data exists before investing in the regression analysis.

Our key finding: HF Hub and OpenML exhibit strong domain and temporal selection biases aligned with their founding communities, not with the broader benchmark ecosystem of the pre-2018 ML literature. HF Hub emerged from the NLP/transformer community --- IMDB, SST-2, GLUE are present; large-scale DL vision (ImageNet, COCO), speech (LibriSpeech, TIMIT), and graph (Cora, Citeseer) datasets are absent. Small-scale image benchmarks widely used in NLP comparisons (CIFAR-10, CIFAR-100, SVHN) are partially present on HF Hub but do not represent the large-scale DL vision datasets that dominate Raff's post-2012 papers. OpenML was designed for classical tabular ML experiments and architecturally cannot host large image or audio datasets --- iris, adult, covertype are present; CIFAR, ResNet benchmarks are absent. This is not a pipeline defect --- 18/18 tests pass and the pipeline runs against live APIs in under 90 seconds. The coverage gap is in the platforms themselves.

This paper makes the following contributions:

\textbf{C1 (Empirical):} The first systematic quantification of HuggingFace Hub and OpenML coverage for the Raff 2019 reproducibility corpus, establishing that HF Hub achieves 30\% dataset card coverage (concentrated in NLP benchmarks and small-scale image classification benchmarks; large-scale DL vision, speech, and graph datasets absent) and OpenML achieves 22\% pre-publication temporal coverage (restricted to classical tabular datasets). Both are insufficient for valid regression analysis.

\textbf{C2 (Infrastructure finding):} Documentation of the domain-temporal selection bias in both APIs: HF Hub's ecosystem reflects its NLP/transformer origin; OpenML's architecture structurally excludes large image, audio, and graph datasets. Any reproducibility auditing study of pre-2018 ML literature using these APIs will be structurally biased toward NLP papers.

\textbf{C3 (Methodological):} Identification of Papers With Code API as the appropriate primary data source for this class of reproducibility auditing study, motivated by our characterization of HF+OpenML limitations.

\textbf{C4 (Reusable infrastructure):} A validated, 18-test data acquisition pipeline (RaffParser, HFCoverageChecker, OpenMLTemporalChecker, MetricsAggregator) that is immediately reusable for a redesigned study using Papers With Code as primary API.

We organize the paper as follows. Section~\ref{sec:related} reviews the dataset documentation and ML reproducibility literatures and positions our infrastructure study within them. Section~\ref{sec:methodology} describes the pipeline design and experimental protocol. Section~\ref{sec:experiments} presents the experimental setup. Section~\ref{sec:results} presents the coverage measurements and root-cause analysis. Section~\ref{sec:discussion} discusses implications and the path to the redesigned study. Section~\ref{sec:conclusion} concludes.
